Fix a transform \(M\) represented in the class.  By \(\text{lem:zero-localization}\), it has a zero in \(|z|\le R_0\), so every circle with radius larger than \(R_0\) encloses at least one zero.  The transform bound (1) in that lemma gives
\[
 \log\max_{|z|=R_{\mathrm{out}}}|M(z)|
 \le4\psi_{\eta}^2R_{\mathrm{out}}^2.
\]
Because \(M(0)=1\), Jensen's formula shows that the number of zeros in \(|z|\le R_{\mathrm{fac}}\), counted with multiplicity, is at most
\[
 Q_{\star}
 =\frac{4\psi_{\eta}^2R_{\mathrm{out}}^2}
        {\log(R_{\mathrm{out}}/R_{\mathrm{fac}})}
 \le N_{\mathrm{cert}}. \tag{1}
\]

The grid \(\mathcal Q_{m_{\star}}\) is finite, lies in \((R_0,R_1)\), contains more than \(N_{\mathrm{cert}}\) points, and has spacing \(d_{\star}\).  A zero modulus can lie within \(h_{\star}=d_{\star}/3\) of at most one grid point.  Hence (1) and the pigeonhole principle supply a grid radius \(\rho_j\) whose circle is at radial distance at least \(h_{\star}\) from every zero in \(|z|\le R_{\mathrm{fac}}\).  This circle encloses the zero already localized in \(|z|\le R_0\), and therefore \(N_{C_j}(M)\ge1\).

Factor the at most \(N_{\mathrm{cert}}\) zeros in the disk by their disk Blaschke factors and denote the remaining zero-free factor, normalized at the origin, by \(U\).  On \(|z|\le R_1\), radial separation bounds the product of the Blaschke moduli from below by
\[
 \left(\frac{h_{\star}}{R_{\mathrm{fac}}+R_1}\right)^{N_{\mathrm{cert}}}.
\]
Applying the Poisson formula to the positive harmonic majorant of \(-\log|U|\), using the transform bound on \(|z|=R_{\mathrm{fac}}\), gives
\[
 \inf_{|z|\le R_1}|U(z)|
 \ge
 \exp\{-4(\Lambda_{\star}-1)\psi_{\eta}^2R_{\mathrm{fac}}^2\}.
\]
Multiplication yields
\[
 \inf_{z\in C_j}|M(z)|\ge a_{\mathrm{an}}\ge a_{\star}>0. \tag{2}
\]

It remains to verify effectivity.  Certified-real names permit terminating outward refinement for the integer bound \(N_{\mathrm{cert}}\ge Q_{\star}\).  The search for \(m_{\star}\), enumeration of its finite rational grid, and all subsequent arithmetic are finite.  Since \(a_{\mathrm{an}}>0\), outward interval refinement also eventually produces a positive dyadic rational \(a_{\star}\le a_{\mathrm{an}}\).  Every input to these operations is a displayed primitive constant; no value of the unknown transform \(M\), and hence no knowledge of the law of \(\eta\), is used.  The same computed bank and certificate work for every \(n\) by (1)--(2).